\documentclass[11pt]{article}

% ========= 基本包 =========
\usepackage[margin=1in]{geometry}
\usepackage[T1]{fontenc}
\usepackage[utf8]{inputenc}
\usepackage{lmodern}
\usepackage{amsmath, amssymb, amsthm, amsfonts}
\usepackage{mathtools}
\usepackage{bm}
\usepackage{enumitem}
\usepackage{graphicx}
\usepackage{xcolor}
\usepackage{hyperref}
\usepackage{cite}
\usepackage{microtype}

% ========= 超链接样式 =========
\hypersetup{
  colorlinks=true,
  linkcolor=blue!60!black,
  citecolor=blue!60!black,
  urlcolor=blue!70!black
}

% ========= 定理环境 =========
\newtheorem{theorem}{Theorem}[section]
\newtheorem{lemma}[theorem]{Lemma}
\newtheorem{proposition}[theorem]{Proposition}
\newtheorem{corollary}[theorem]{Corollary}

\theoremstyle{definition}
\newtheorem{definition}[theorem]{Definition}
\newtheorem{assumption}[theorem]{Assumption}
\newtheorem{example}[theorem]{Example}

\theoremstyle{remark}
\newtheorem{remark}[theorem]{Remark}

% ========= 常用宏 =========

% 向量和范数
\newcommand{\R}{\mathbb{R}}
\newcommand{\Rn}{\mathbb{R}^n}
\newcommand{\N}{\mathbb{N}}
\newcommand{\inner}[2]{\left\langle #1, #2 \right\rangle}
\newcommand{\norm}[1]{\left\lVert #1 \right\rVert}
\newcommand{\abs}[1]{\left\lvert #1 \right\rvert}

% 正锥、Hilbert 度量等
\newcommand{\cone}{\mathcal{K}}
\newcommand{\Pcone}{\mathbb{P}(\cone)}
\newcommand{\Hilbert}[2]{d_H\!\left(#1,#2\right)}
\newcommand{\hilbert}{d_H}
\newcommand{\proj}{\mathbb{P}}
\newcommand{\Id}{\mathrm{Id}}

% 优化、梯度流
\newcommand{\grad}{\nabla}
\newcommand{\dd}{\,\mathrm{d}}
\newcommand{\dt}{\,\mathrm{d}t}

% 备注环境（写给自己看的）
\newcommand{\todo}[1]{\textcolor{red}{[TODO: #1]}}

% ========= 文档信息 =========

\title{A Preliminary Note on Hilbert's Projective Geometry in Gradient-Based Learning Dynamics}
% 如果你想在标题下面写 “Preliminary note for Francesco Tudisco”，可以加一行小字：
% \title{A Preliminary Note on Hilbert's Projective Geometry in Gradient-Based Learning Dynamics\\[0.3em]
% \large Preliminary note for Francesco Tudisco}

\author{Xinyang (Elizabeth) Wen\\
\small University of Wisconsin--Madison\\
\small \texttt{xwen57@wisc.edu}
}

\date{\today}

% ========= 正文开始 =========
\begin{document}

\maketitle

\begin{abstract}
This short note collects some preliminary observations on gradient-based learning dynamics viewed through the lens of Hilbert's projective metric on a positive cone. 
In a simple gradient-descent setting with a nonnegative parametrization, we observe that the iterates exhibit a Perron--Frobenius--type behaviour: for most of training, the projective distance to the final parameter vector remains almost constant, while a sudden contraction appears in the late phase of optimization. 
We briefly describe the experimental setup, the empirical phenomena, and some speculative connections with nonlinear Perron--Frobenius theory and cone contractions.
\end{abstract}

% ========== 1. Introduction ==========
\section{Introduction}

% 这里写动机和背景，尽量简洁。
Gradient-based methods are the workhorse of modern machine learning, yet a detailed geometric description of their trajectories remains challenging, especially in high-dimensional and nonconvex settings.
In this note, we explore a simple but, to our knowledge, nonstandard viewpoint: we treat the parameter trajectory as living in a positive cone and measure its evolution using Hilbert's projective metric.

Very roughly, in models with a nonnegative parametrization (induced, for instance, by an appropriate non-linear activation and loss choice), the parameter iterates remain in a cone $\cone \subset \Rn_{>0}$ and appear to exhibit a Perron--Frobenius--type contraction in the Hilbert projective metric.
Empirically, we observe that the Hilbert distance to the final parameter vector is nearly constant throughout most of training and collapses rapidly only in a short late phase.

The purpose of this note is \emph{not} to provide a full theory, but rather:
\begin{itemize}[leftmargin=*]
    \item to specify a minimal setup where the cone structure and Hilbert geometry arise naturally;
    \item to document the empirical contraction phenomena we observe in simple experiments (e.g., on MNIST);
    \item to sketch some speculative connections with nonlinear Perron--Frobenius theory and Birkhoff-type contractions.
\end{itemize}

\paragraph{Acknowledgment.}
This note is written in response to a kind email exchange with Francesco Tudisco about possible links between gradient flows and Perron--Frobenius theory.

% ========== 2. Setup and notation ==========
\section{Setup and notation}

In this section we outline a simple gradient-descent model with a nonnegative parametrization and fix some notation.

\subsection{Positive parametrization}

% 在这里写下你具体的模型假设，可以慢慢补。
We consider a parameter vector $w_t \in \Rn$ evolving under gradient descent on a loss $L:\Rn \to \R$:
\begin{equation}
    w_{t+1} = w_t - \eta \grad L(w_t), \qquad t = 0,1,2,\dots,
    \label{eq:GD}
\end{equation}
with a fixed step-size $\eta > 0$.
We assume that the parametrization is such that the effective parameters remain strictly positive; for example, one can think of a map of the form
\[
    w_t = \phi(\theta_t),
\]
where $\theta_t$ are the raw weights and $\phi$ is a componentwise nonlinearity with $\phi(\R) \subset (0,\infty)$ (e.g., Softplus or a shifted exponential).

\begin{assumption}[Positive cone invariance]
\label{assumption:cone}
There is a closed convex cone $\cone \subset \Rn$ with nonempty interior such that $w_t \in \mathrm{int}(\cone)$ for all $t \ge 0$.
\end{assumption}

In our experiments, $\cone$ is simply the standard positive cone $\R^n_{>0}$, arising from elementwise positive parametrizations of the last layer of a neural network. 
We focus on the dynamics of $w_t$ within this cone.

\subsection{Gradient flow limit}

For conceptual clarity, it is sometimes convenient to consider the continuous-time gradient flow
\begin{equation}
    \frac{\dd}{\dt} w(t) = - \grad L(w(t)), \qquad w(0) \in \mathrm{int}(\cone),
    \label{eq:GF}
\end{equation}
and view \eqref{eq:GD} as an Euler discretization.
Our empirical observations, however, are directly based on discrete-time iterates.

% ========== 3. Hilbert projective metric ==========
\section{Hilbert's projective metric on a positive cone}

We briefly recall Hilbert's projective metric and fix notation.
We refer to standard references on cone metrics and Perron--Frobenius theory for background.

\begin{definition}[Hilbert projective metric]
Let $\cone \subset \Rn$ be a proper cone and let $x,y \in \mathrm{int}(\cone)$.
Define
\[
    M(x,y) := \inf \{ \alpha > 0 : x \le \alpha y \}, \qquad 
    m(x,y) := \sup \{ \beta > 0 : \beta y \le x \},
\]
where inequalities are with respect to the partial order induced by $\cone$.
The \emph{Hilbert projective metric} is
\[
    \Hilbert{x}{y} := \log \frac{M(x,y)}{m(x,y)}.
\]
\end{definition}

In the special case of the standard cone $\R^n_{>0}$, one has the familiar expression
\begin{equation}
    \Hilbert{x}{y}
    = \log \left(
        \frac{\max_i x_i / y_i}{\min_i x_i / y_i}
    \right),
    \qquad x,y \in \R^n_{>0}.
    \label{eq:hilbert-positive-orthant}
\end{equation}

\begin{remark}[Projective nature]
The metric $\hilbert$ is invariant under positive scaling:
$\Hilbert{\lambda x}{\mu y} = \Hilbert{x}{y}$ for all $\lambda,\mu>0$.
Thus it is naturally a metric on the projective space $\proj(\cone)$ of rays in $\cone$.
\end{remark}

\begin{remark}[Positive maps and contractions]
Classical results (Birkhoff, Bushell, and many others) show that order-preserving, homogeneous maps $T:\cone \to \cone$ with bounded projective diameter are nonexpansive or even contractive in $(\proj(\cone),\hilbert)$.
Our empirical observations suggest that late-phase gradient iterates behave \emph{approximately} in this spirit, even though the gradient map is not exactly homogeneous.
\end{remark}

% ========== 4. Empirical observations ==========
\section{Empirical observations in gradient-based learning}

In this section we summarize the empirical phenomena observed in simple experiments. 
The precise experimental details can be refined or expanded as needed.

\subsection{Experimental setup (sketch)}

% TODO: 在这里写你实际用的模型、数据、参数等
\todo{Describe the specific model (e.g., a simple fully connected network on MNIST, with a positive last layer), the loss function, optimization hyperparameters, and how $w_t$ is extracted from the network.}

For each run, we:
\begin{enumerate}[leftmargin=*]
    \item Initialize $w_0 \in \R^n_{>0}$ (e.g., from a standard initialization mapped through a positive nonlinearity).
    \item Train for a fixed number of steps $T$ using gradient descent or a closely related optimizer.
    \item Save the trajectory $\{w_t\}_{t=0}^T$ of the last-layer parameters.
    \item Compute various Hilbert distances and related quantities along the trajectory.
\end{enumerate}

\subsection{Hilbert distance to the final parameter}

Let $w^\ast := w_T$ denote the final parameter vector.
We measure the Hilbert distance to the final point:
\[
    h_t := \Hilbert{w_t}{w^\ast}, \qquad t = 0,1,\dots,T.
\]

\paragraph{Qualitative behaviour.}
Empirically, we observe the following striking pattern:
\begin{itemize}[leftmargin=*]
    \item For a large portion of the training trajectory, the sequence $(h_t)$ is almost flat:
    \[
        h_{t+1} \approx h_t, \quad \frac{h_{t+1}}{h_t} \approx 1,
    \]
    indicating that the \emph{projective direction} of $w_t$ relative to $w^\ast$ remains essentially unchanged.
    \item In a relatively short final phase, $h_t$ collapses rapidly to $0$, corresponding to $w_t$ aligning with the ray spanned by $w^\ast$ in the cone.
\end{itemize}

% TODO: 你可以在这里插入图
% \begin{figure}[h]
%     \centering
%     \includegraphics[width=0.6\textwidth]{hilbert_to_final_example.pdf}
%     \caption{Example of Hilbert distance to the final parameter $w^\ast$ along training.}
%     \label{fig:hilbert-to-final}
% \end{figure}

\subsection{Between-step distances and direction freezing}

We also consider the between-step Hilbert distances
\[
    b_t := \Hilbert{w_t}{w_{t+1}},
\]
and Euclidean or cosine distances for comparison.

% T
